\section{The complete ordinary localization map}\label{sec:localization}
Let $z$ be an ordinary $k$ valued character of $\Gsig$. It is a difference of bare primitives, not beta and not either primitive itself. Local class field theory gives the following coordinate order:
\begin{equation}\label{eq:localcoords}
F_2,F_3,F_5,\quad7:\phi,3,7,\quad
I_{q_1},F_{q_1},\ldots,I_{q_9},F_{q_9}.
\end{equation}
We use the local Artin convention sending a selected uniformizer to arithmetic Frobenius. At five the uniformizer is $\sqrt5$. At support primes it is the rational prime $q$. Every sign in the reciprocity rows and every value of $h=(Nv-1)/7$ below uses this convention \cite{milneCFT2020}.

At $2,3,5$, the residue norms are $4,9,5$, whose unit orders have no seven factor. Only the unramified coordinate remains. At seven, the two independent pro seven unit coordinates are $\phi,3$, whose normalized logarithmic first digits are $5\sqrt5,6$, and the third coordinate is the uniformizer seven. Each of the nine inert support primes has one seventh power residue coordinate and one unramified coordinate. Real odd primary characters vanish. The total local dimension is
\[
3+3+2\cdot9=24\quad\text{over }k.
\]
The chosen $g_q=\phi+a_q$ generating the residue seventh power quotient supplies $I_q$. It need not generate the whole multiplicative group. Only its nonzero seventh power symbol and specified normalization are required.

\subsection{Fourteen principal relations are exhaustive}
The $\Sigma$ unit group modulo seventh powers has basis
\begin{equation}\label{eq:sigmabasis}
\phi,2,3,\sqrt5,7,q_1,\ldots,q_9.
\end{equation}
Prime valuations prove independence of the thirteen prime coordinates, and \eqref{eq:phi-nonseventh} proves the unit coordinate is nontrivial. Class number one removes extra ideal class variables. Principal reciprocity pairs these fourteen units with the twenty four local characters and gives a matrix $L$ of size $14$ by $24$.

There are two complementary reasons why this matrix has rank fourteen. Lemma~\ref{lem:sha1} proves that global Kummer localization is injective. Under the local perfect pairings, $L$ is the transpose of the global Kummer localization matrix, so the fourteen global Kummer coordinates give fourteen independent functionals. The finite arithmetic then reconstructs their declared coordinate entries and verifies this rank with explicit certificates. This order avoids treating a computational rank as the only reason a global reciprocity presentation is complete.

The finite set Poitou--Tate sequence identifies the entire ordinary localization image with $\ker L$. Equivalently, the ray class presentation says that a compatible family of local characters extends exactly when it annihilates the principal $\Sigma$ units \cite{milneADT2006,milneCFT2020}. This is a restricted image theorem. It is not unrestricted finite localization surjectivity.

Ordinary global localization is injective here. A character with zero local classes at every allowed ramification place is everywhere unramified. Odd real terms vanish, and the ordinary class group is trivial. Thus it is zero. The ordinary global space is therefore ten dimensional, since $24-14=10$.

\subsection{An explicit image parametrization}
Put $a=z_7(3)$ and $x_i=z(I_{q_i})$. Let $\ell_i$ be the seventh power logarithm normalized by $\ell_i(g_{q_i})=1$. The supporting symbols of $\phi$ vanish. For rational seven adic units define
\[
\ell_7(t)=\frac{\log(t^6)}{\log(3^6)}\pmod7
=6\frac{t^6-1}{7}\pmod7.
\]
Principal reciprocity determines the remaining coordinates:
\begin{align}
z_7(\phi)&=0,\label{eq:Jphi}\\
z(F_2)&=-5a-\sum_i\ell_i(2)x_i,\label{eq:J2}\\
z(F_3)&=-a-\sum_i\ell_i(3)x_i,\label{eq:J3}\\
z(F_5)&=-4a-\sum_i\ell_i(\sqrt5)x_i,\label{eq:J5}\\
z_7(7)&=-\sum_i\ell_i(7)x_i,\label{eq:J7}\\
z(F_q)&=-\ell_7(q)a-\sum_{i:q_i\ne q}\ell_i(q)x_i.\label{eq:Jq}
\end{align}
At its own conductor place, the rational prime $q$ is the chosen uniformizer with unit part one. Its self residue logarithm in the ordinary character presentation is zero, not the undefined logarithm of zero. Every other support prime contributes simultaneously.

These equations define a $24$ by $10$ matrix $J$. The ten free parameters occur in their declared local coordinates, so $J$ is injective. Direct substitution gives $LJ=0$. Since both $\im J$ and $\ker L$ have dimension ten,
\begin{equation}\label{eq:Jcomplete}
\im J=\ker L.
\end{equation}
Every parameter vector gives a global ordinary character. This is not merely an overapproximation by necessary equations.

The matrices have entries in $\F_7$ and are extended to $k$. Their use does not assume coefficient descent. A smaller common pointwise semilinear problem can be imposed afterward, but it requires the same $j$, frame, and reference for the compared cochains.
